\documentclass[10pt,a4paper]{amsart}
\usepackage[margin=19mm]{geometry}
\usepackage{amsmath,amssymb,mathtools}
\usepackage[scr=rsfs,cal=euler]{mathalfa}
\usepackage{xfrac}
\usepackage[unicode,hidelinks,bookmarks=false]{hyperref}
\newcommand{\ie}{i.e.\ }
\newcommand{\cf}{cf.\ }
\newcommand{\resp}{respectively\ }
\newcommand{\Subsection}{\subsection}
\providecommand{\nobreakdash}{\nobreak\hbox{-}}
\setlength{\emergencystretch}{2em}
\multlinegap0pt
\usepackage{amssymb}
\usepackage{enumerate}
\usepackage{algorithm}\usepackage{algpseudocode}
\usepackage{multirow}
\newtheorem{lemma}{Lemma}
\title{Computational Complexity of Finding Subgroups of a Given Order}
\author{K. Lakshmanan}
\date{Source: arXiv:2503.11238v2 (CC BY 4.0)}
\begin{document}
\maketitle

\noindent\emph{Source and licence.} Computational Complexity of Finding Subgroups of a Given Order. Version arXiv:2503.11238v2 (CC BY 4.0), distributed by arXiv under the Creative Commons Attribution 4.0 International licence, http://creativecommons.org/licenses/by/4.0/, as recorded in the arXiv licence metadata for this exact version and shown on the abstract page https://arxiv.org/abs/2503.11238v2. Full licence notice: \texttt{licenses/arxiv-2503.11238-CC-BY-4.0.txt}.\par\smallskip

\noindent\textit{Editorial scope.} Counted unit: the article's lemma generatedequals (source0.tex lines 201--206). The article's complete original proof of it is delivered with it (source0.tex lines 208--215, one \texttt{proof} environment of 8 lines). No mathematics is written, repaired or re-derived: the statement and the proof are the article's own lines, in the article's own notation, and no retained line is substituted or re-flowed.  Retained uncounted as the source-local context the proof needs: lines 183--185.  Nothing was omitted from any retained span, so the package carries no omission record.  No retained line contains a citation, so the wrapper carries no bibliography and no bibliography entry.  No retained line contains a figure environment, an image-inclusion command or drawing code, so no image is delivered and the package declares no asset dependency.  Editorial formatting only, in addition: an AMS \texttt{amsart} preamble that restates the article's own declarations and macros used by the retained lines, namely the environments \texttt{lemma} on separate counters, together with the standard packages those lines need.  The retained environments print in this document as Lemma 1 in order of appearance, whereas the article prints the same items with the numbering of its own full document.  The complete native source of the article as distributed in the arXiv e-print for this exact version is bundled unchanged as \texttt{sources/174747/source0.tex}.
	\begin{lemma}\label{coprimelemma}
		Let $a \in G$ satisfy $\gcd(|a|,m)=1$. Then no subgroup of $G$ of order $m$ contains any nontrivial element of $\langle a \rangle$.
	\end{lemma}

	\begin{lemma}\label{generatedequals}
		Let $a_1,\dots,a_i$ be the elements retained by the algorithm. Then
		\[
		\langle a_1,\dots,a_i\rangle \;=\; G_m.
		\]
	\end{lemma}

	\begin{proof}
		By Lemma~\ref{coprimelemma}, only elements whose order shares a prime divisor with $m$ are retained. Hence each selected element lies in $G_m$, so
		\[
		\langle a_1,\dots,a_i\rangle \subseteq G_m.
		\]
		
		Conversely, suppose some element of $G_m$ were not contained in $\langle a_1,\dots,a_i\rangle$. Then some nontrivial $p$-component for $p \mid m$ would remain outside the generated subgroup. Since the algorithm processes all elements not removed, such elements would eventually be selected, contradicting termination. Hence equality holds.
	\end{proof}

\end{document}
